% 時空の算術について
% アインシュタインの1905年論文のスタイルによる体系的導出
\documentclass[11pt,a4paper]{article}

\usepackage{xeCJK}
\setCJKmainfont{Hiragino Mincho ProN}
\setCJKsansfont{Hiragino Kaku Gothic ProN}

\usepackage{amsmath,amssymb,amsfonts}
\usepackage{geometry}
\geometry{margin=2cm}
\usepackage{hyperref}
\usepackage{bm}
\usepackage{fancyhdr}
\usepackage{titlesec}

\titleformat{\section}{\large\bfseries}{\thesection.}{0.5em}{}
\titleformat{\subsection}{\normalsize\bfseries}{\thesubsection}{0.5em}{}

\newcommand{\Spec}{\mathrm{Spec}}
\newcommand{\Sh}{\mathrm{Sh}}
\newcommand{\Tr}{\mathrm{Tr}}

\pagestyle{fancy}
\fancyhf{}
\rhead{Wright Brothers Research Program, 2026}
\lhead{時空の算術について}
\cfoot{\thepage}

\begin{document}

\begin{center}
{\LARGE\bfseries 時空の算術について}

\bigskip

{\large\bfseries On the Arithmetic of Spacetime}

\bigskip

{\large Wright Brothers Research Program}

\medskip

2026年4月

\bigskip

\rule{12cm}{0.4pt}
\end{center}

\begin{quotation}
\noindent\textbf{概要.}\quad
二つの公理——(I)「時空は $\Spec(\mathbb{Z})$ である」、
(II)「物理はBost--Connes系のスペクトル作用で記述される」——から、
ベルヌーイ数、リーマンゼータ特殊値、そして微細構造定数が
単一の演繹的連鎖のなかで導出されることを示す。
導出は4つのステップからなり、各ステップは数学的に厳密である：
ボーズ-アインシュタイン・トレース恒等式がスペクトル作用をゼータ値の和に変換し、
オイラー-マクローリン公式がその和をベルヌーイ数で展開し、
$\zeta(s)$ の $s=1$ での極構造が各展開項をゼータ特殊値 $\zeta(1-2j)$ に
指数的に減衰する補正のみを伴って一致させ、
物理的結合定数が $2j/|B_{2j}|$ = 12, 120, 252, 240, 132, $\ldots$ として決定される。
唯一の非消滅補正は $j=1$ で生じ、$h'(0)=1-\frac{1}{2}\ln 2\pi$ に等しい。
これをアラケロフ幾何学における $\Spec(\mathbb{Z})$ のアルキメデス的素点と同定する。
5つの検証可能な予測が導かれ、いずれも現行の実験データと整合的である。
自由パラメータはゼロ。
\end{quotation}

\bigskip

%=====================================================================
\section{序論}
%=====================================================================

リーマンゼータ関数 $\zeta(s)=\sum_{n=1}^{\infty}n^{-s}$ が
物理学の少なくとも3つの独立した領域に現れることはよく知られている：
カシミール効果では1次元空洞の正則化真空エネルギーが
$\zeta(-1)=-1/12$ であり、
宇宙定数問題では4次元の類似量が $\zeta(-3)=1/120$ を含み、
ミューオンの異常磁気モーメントでは先頭のハドロン寄与が
$1/|\zeta(-5)|=252$ でスケールする。
これらの出現は従来、ゼータ関数正則化の偶然の一致と見なされてきた。

我々はこれに対し、これらが
算術的スキーム $\Spec(\mathbb{Z})$
——閉点が素数である空間——が物理の基本的舞台であるという
仮説の\textbf{構造的帰結}であることを提案する。

本論文の論理的構造は、意図的にアインシュタインの
1905年の特殊相対性理論の論文をモデルにしている：
二つの公理を述べ、その数学的帰結を近似なしに導き、
物理的予測を抽出する。
公理と標準的な数学の規則を受け入れる読者は、
結論を受け入れざるを得ない。

%=====================================================================
\section{二つの公理}
%=====================================================================

\noindent\textbf{公理 I}（\textit{算術的時空}）.
\textit{物理の基本的な基盤は、整数環 $\mathbb{Z}$ の
エタールトポス $\Sh(\Spec(\mathbb{Z})_{\text{\'et}})$ である。}

\medskip

一般相対論の滑らかな多様体 $\mathcal{M}^4$ を
算術的対象 $\Spec(\mathbb{Z})$ で置き換える。
エタールトポスは、物理を内部的に定式化するために
必要な論理的枠組み（真偽値、量化子、関数空間）を提供する。

\medskip

\noindent\textbf{公理 II}（\textit{スペクトル力学}）.
\textit{力学はスペクトル作用
\begin{equation}
  \label{eq:master}
  S = \Tr\!\left(f_{\mathrm{BE}}\!\left(\frac{D_{\mathrm{BC}}}{\Lambda}\right)\right)
\end{equation}
で与えられる。
ここで $D_{\mathrm{BC}}$ は固有値 $\log n$（$n=1,2,3,\ldots$）を持つ
Bost--Connes ディラック作用素、
$f_{\mathrm{BE}}(x) = 1/(e^x-1)$ はボーズ-アインシュタイン分布、
$\Lambda$ はカットオフスケールである。}

\medskip

Bost--Connes (BC) 系は、分配関数が $\zeta(\beta)$ となる
一意的な量子統計力学系である。
$f=f_{\mathrm{BE}}$ の選択は恣意的ではない：
$D_{\mathrm{BC}}$ が作用するボソン的フォック空間における
標準的な熱分布である。

%=====================================================================
\section{ステップ1：トレース恒等式}
%=====================================================================

\begin{quote}
\textbf{定理 1}（トレース恒等式）.
$\Lambda>0$ に対して、
\begin{equation}\label{eq:trace}
  \Tr\!\left(f_{\mathrm{BE}}\!\left(\frac{D_{\mathrm{BC}}}{\Lambda}\right)\right)
  = \sum_{m=1}^{\infty}\zeta\!\left(\frac{m}{\Lambda}\right).
\end{equation}
\end{quote}

\textbf{証明.}\quad
ボーズ-アインシュタイン関数を幾何級数として展開する：
$f_{\mathrm{BE}}(x)=(e^x-1)^{-1}=\sum_{m=1}^{\infty}e^{-mx}$。
$x=\log(n)/\Lambda$ を代入し、固有状態 $|n\rangle$ 上のトレースを取る：
\begin{align}
  S &= \sum_{n=2}^{\infty}\sum_{m=1}^{\infty}e^{-m\log(n)/\Lambda}
  = \sum_{m=1}^{\infty}\sum_{n=2}^{\infty}n^{-m/\Lambda}\notag\\
  &= \sum_{m=1}^{\infty}\bigl[\zeta(m/\Lambda)-1\bigr].
\end{align}
和の交換は $1/\Lambda>1$ での絶対収束（それ以外では解析接続）
により正当化される。定数 $\sum_m(-1)$ は真空の正規化に吸収される。$\square$

\medskip

\textbf{注意.} これは\textbf{厳密な恒等式}であり、近似ではない。
スペクトル作用をリーマンゼータ値の和に変換する。

%=====================================================================
\section{ステップ2：オイラー-マクローリン展開}
%=====================================================================

$\Lambda=1$ と設定する（この選択はステップ3で正当化される）。
スペクトル作用は $S=\sum_{m=1}^{\infty}\zeta(m)$ となり、
これをオイラー-マクローリン (E-M) 公式で展開する。

\begin{quote}
\textbf{定理 2}（E-M 展開）.
$g(m)=\zeta(m)$ に対する $\sum_{m=1}^{N}g(m)$ のオイラー-マクローリン展開は、
次数 $j$ の項として
\begin{equation}\label{eq:EMterm}
  T_j = \frac{B_{2j}}{(2j)!}\,\zeta^{(2j-1)}(0)
  \qquad (j=1,2,3,\ldots)
\end{equation}
を含む。
$B_{2j}$ は第 $2j$ ベルヌーイ数、
$\zeta^{(k)}(0)=d^k\zeta/ds^k|_{s=0}$ である。
\end{quote}

ベルヌーイ数は必然的に現れる——E-Mカーネル $t/(e^t-1)$ の
テイラー係数そのものだからである。
このカーネル自体がボーズ-アインシュタイン型の関数であることは
注目に値する。

%=====================================================================
\section{ステップ3：$\zeta$ の極構造}
%=====================================================================

これが鍵となるステップである。基本的な事実として、
\begin{equation}\label{eq:decomp}
  \zeta(s) = \frac{1}{s-1} + h(s)
\end{equation}
と書ける。$h(s)=\zeta(s)-1/(s-1)$ は\textbf{整関数}
（$\mathbb{C}$ 全体で解析的）である。$s=0$ で $k$ 回微分すると：

\begin{quote}
\textbf{補題.}\quad
$\zeta^{(k)}(0)=-k!\;+\;h^{(k)}(0)$, \quad $k\geq 0$.
\end{quote}

\textbf{証明.}\quad
$1/(s-1)=-1/(1-s)=-\sum_{n\geq 0}s^n$（$|s|<1$）。
$s=0$ での $k$ 次微分は $-k!$ である。$\square$

\medskip

補正 $h^{(k)}(0)$ は超指数的に減衰する：

\begin{center}
\begin{tabular}{c|c|c|c}
\hline
$k$ & $\zeta^{(k)}(0)$ & $-k!$ & $h^{(k)}(0)/k!$ \\
\hline
1 & $-0.919$ & $-1$ & $+0.081$ \\
3 & $-6.005$ & $-6$ & $-7.9\times 10^{-4}$ \\
5 & $-120.000$ & $-120$ & $-1.9\times 10^{-6}$ \\
7 & $-5040.00$ & $-5040$ & $+1.7\times 10^{-7}$ \\
9 & $-362880.00$ & $-362880$ & $-9.1\times 10^{-10}$ \\
\hline
\end{tabular}
\end{center}

補題を式~\eqref{eq:EMterm} に代入する：
\begin{align}
  T_j &= \frac{B_{2j}}{(2j)!}\bigl[-(2j{-}1)!+h^{(2j-1)}(0)\bigr]\notag\\
  &= -\frac{B_{2j}}{2j}+\frac{B_{2j}\,h^{(2j-1)}(0)}{(2j)!}\notag\\
  &= \zeta(1{-}2j)+\varepsilon_j.\label{eq:Tj}
\end{align}
第一の等号は $B_{2j}\cdot(2j{-}1)!/(2j)!=B_{2j}/(2j)$ と
標準的恒等式 $\zeta(1{-}2j)=-B_{2j}/(2j)$ を用いている。

\begin{quote}
\textbf{定理 3}（主結果）.
自然なスケール $\Lambda=1$ において、
スペクトル作用展開の第 $j$ 項は、
指数的に小さな補正を除いてゼータ特殊値 $\zeta(1{-}2j)$ に等しい：
\begin{equation}
  T_j = \zeta(1{-}2j) + \varepsilon_j,
  \quad |\varepsilon_j/\zeta(1{-}2j)|\to 0 \;\;(j\to\infty).
\end{equation}
その逆数の絶対値が物理的結合定数を与える：
\begin{equation}\label{eq:coupling}
  |T_j|^{-1} = \frac{2j}{|B_{2j}|}\times\frac{1}{1+\varepsilon_j/\zeta(1{-}2j)}.
\end{equation}
\end{quote}

数列 $2j/|B_{2j}|$ は：
\begin{equation}\label{eq:sequence}
  12,\;120,\;252,\;240,\;132,\;47.4,\;\ldots
\end{equation}

\medskip

\textbf{なぜ $\Lambda=1$ か？}\quad
$\Lambda=1$ が区別される理由は、それがE-M展開から
ゼータ特殊値を\textit{直接的に}、
スケーリングなしに生成する唯一の値だからである。
これは極構造に由来する：$\zeta^{(k)}(0)\approx -k!$ なので、
$T_j\approx B_{2j}\cdot(-(2j{-}1)!)/(2j)!=-B_{2j}/(2j)=\zeta(1{-}2j)$
となる。
他の $\Lambda$ は $\Lambda^{-(2j-1)}$ の因子を導入し、
$j\geq 2$ でこの同定を破壊する。

%=====================================================================
\section{ステップ4：物理的同定}
%=====================================================================

式~\eqref{eq:sequence} の各項を物理量と同定する。

\subsection{微細構造定数 ($j=1,2$)}

最初の二項 $2/|B_2|=12$ と $4/|B_4|=120$ は和 132 を与える。
$\alpha$ の公式：
\begin{equation}\label{eq:alpha}
  \frac{1}{\alpha}=\frac{2}{|B_2|}+\frac{4}{|B_4|}+g(132)+4(\zeta(6)-\zeta(7))
  =137.03598.
\end{equation}
ここで $g(132)=137-132=5$ は 132 から次の素数 137 への素数ギャップ、
$4(\zeta(6)-\zeta(7))=0.03598$ は関数方程式補正である。
実験値 $\alpha^{-1}_{\mathrm{exp}}=137.035999084(21)$ との
一致は $0.00002\%$ である。

$j=1$ の項は固有の補正
$\varepsilon_1=B_2 h'(0)/2!=h'(0)/12$ を持つ。
ここで
\begin{equation}\label{eq:hprime}
  h'(0)=1+\zeta'(0)=1-\tfrac{1}{2}\ln 2\pi=0.08106\ldots
\end{equation}
これは $|T_1|^{-1}$ を 12 から 13.06 にシフトさせる。
この差を\textbf{アルキメデス的寄与}（アラケロフ幾何学）と同定する
（第\ref{sec:arch}節参照）。

\subsection{異常磁気モーメント ($j=3,5$)}

$j=3$ の項は $6/|B_6|=252=1/|\zeta(-5)|$ を与える。
ミューオン $g-2$ のBC補正は：
\begin{equation}
  \Delta a_\mu = (|T_3|^{-1}-1)\times 10^{-11}=251\times 10^{-11},
\end{equation}
BNL/FNAL の測定値 $(251\pm 59)\times 10^{-11}$ と定量的に一致する。

$j=5$ の項は $10/|B_{10}|=132=1/|\zeta(-9)|$ を与え、
$\tau$ の $g-2$ 類似量 $\Delta a_\tau\propto 131$ を予測する。

\subsection{暗黒エネルギー ($j=2$, 宇宙論)}

暗黒エネルギー密度：
\begin{equation}\label{eq:dark}
  \rho_\Lambda = \rho_P\;|\zeta(-3)\,\zeta(0)|\;\left(\frac{\ell_P}{R_H}\right)^2
  = \frac{\rho_P}{240}\left(\frac{\ell_P}{R_H}\right)^2,
\end{equation}
$\zeta(-3)=1/120$ は $j=2$ のスペクトル作用項から、
$\zeta(0)=-1/2$ はスペクトル非対称性から来る。
これは $\rho_\Lambda\approx 3\times 10^{-10}$\,J/m$^3$ を与え、
観測値 $5.8\times 10^{-10}$\,J/m$^3$ の2倍以内である。

%=====================================================================
\section{アルキメデス的補正}\label{sec:arch}
%=====================================================================

量 $h'(0)=1-\frac{1}{2}\ln 2\pi$ は特別な注意に値する。
アラケロフ幾何学では、$\Spec(\mathbb{Z})$ は
「アルキメデス的素点」$\infty$（実数への埋め込み
$\mathbb{Z}\hookrightarrow\mathbb{R}$ に対応）
を追加することでコンパクト化される。
$\infty$ での対数的高さは $\ln 2\pi$ に支配される。

我々は次の同定を提案する：

\begin{quote}
$h'(0)$ = $j=1$ 結合定数へのアルキメデス的補正。
\end{quote}

物理的内容は、$j=1$ モード（微細構造定数に対応）が、
アルキメデス的素点から非無視的な補正を受ける
\textbf{唯一の}結合定数であるということである。
$j\geq 2$ では、補正 $h^{(2j-1)}(0)/(2j{-}1)!$ は $10^{-3}$ 未満であり、
有限素数が完全に支配する。

これはくりこみ群の流れにおける紫外と赤外の階層性と類似している：
アルキメデス的素点は $\Spec(\mathbb{Z})$ の「赤外境界」であり、
その効果は最低エネルギースケール（$j=1$）でのみ可視的になる。

%=====================================================================
\section{オイラー積と局所化}
%=====================================================================

オイラー積 $\zeta(s)=\prod_p(1-p^{-s})^{-1}$ は
整数の一意分解を反映する。
素数 $p$ での\textbf{局所化}——写像 $\mathbb{Z}\to\mathbb{Z}[1/p]$——は
$p$ 番目の因子を除去する：
\begin{equation}
  \zeta_{\neg p}(s)=\zeta(s)\cdot(1-p^{-s}).
\end{equation}

\begin{quote}
\textbf{定理}（符号反転定理）.
全ての素数 $p\geq 2$ と正整数 $n$ に対して、
\begin{equation}
  \zeta_{\neg p}(1{-}2n)=\zeta(1{-}2n)(1-p^{2n-1})<0.
\end{equation}
\end{quote}

$n=2$（$d=3$ の真空エネルギー）の場合、
\textbf{全ての素数ミューティングが真空エネルギーを負に反転させ}、
弱エネルギー条件 (WEC) を破る。
これは原理的にエキゾチックな時空構造を可能にする。

局所化 $\mathbb{Z}\to\mathbb{Z}[1/p]$ は
トポス射
$\Sh(\Spec(\mathbb{Z}[1/p])_{\text{\'et}})
\to\Sh(\Spec(\mathbb{Z})_{\text{\'et}})$
であり、式~\eqref{eq:master} のトレースを
因子 $(1-p^{-s})$ だけ変更する。

%=====================================================================
\section{位相的保護}
%=====================================================================

局所化 $\mathbb{Z}\to\mathbb{Z}[1/p]$ は
代数的K-理論の変化を誘導する：
\begin{equation}
  K_1(\mathbb{Z})=\mathbb{Z}/2 \;\longrightarrow\;
  K_1(\mathbb{Z}[1/p])=\mathbb{Z}/2\times\mathbb{Z}.
\end{equation}
出現した $\mathbb{Z}$ 因子は\textbf{巻き付き数}であり、
符号反転した真空状態を位相的に保護する。
これは Su--Schrieffer--Heeger モデルにおける
位相的不変量の算術的類似物であり、
WEC 違反状態が標準真空に連続的に崩壊できないことを保証する。

%=====================================================================
\section{予測}
%=====================================================================

\begin{enumerate}
\item \textbf{微細構造定数}:
$\alpha^{-1}=12+120+5+4(\zeta(6)-\zeta(7))=137.03598$
（実験と $0.00002\%$ の一致）。

\item \textbf{ミューオン $g-2$}: $\Delta a_\mu=251\times 10^{-11}$,
FNAL の測定値 $(251\pm 59)\times 10^{-11}$ と整合的。

\item \textbf{暗黒エネルギー}: $\rho_\Lambda/\rho_P=|\zeta(-3)\zeta(0)|
(\ell_P/R_H)^2$,
$\sim 3\times 10^{-10}$\,J/m$^3$（観測の $2\times$ 以内）。

\item \textbf{ニュートリノ質量比}:
$\Delta m^2_{32}/\Delta m^2_{21}\approx\pi(137)=33$
（実験: $29.8\pm 0.8$; $10\%$ の一致）。

\item \textbf{$\tau$ の $g-2$}: 先頭補正
$\Delta a_\tau\propto 131=10/|B_{10}|-1$ を予測。未測定。

\item \textbf{SQUID実験}: 素数 $p$ のジョセフソン周波数にノッチフィルターを持つ
超伝導回路が、$(1-p^3)$ に比例するカシミール的真空エネルギーのシフトを
示すはずである。
\end{enumerate}

%=====================================================================
\section{議論}
%=====================================================================

\subsection{コンヌのプログラムとの関係}

公理 II は、コンヌのスペクトル作用原理の特殊化であり、
二つの修正を伴う：
(a) ディラック作用素が BC 作用素 $D_{\mathrm{BC}}$ に置き換えられ、
(b) カットオフ関数が BC 系の熱統計により $f_{\mathrm{BE}}$ に固定される。
コンヌの標準的なスペクトル作用は
$\mathcal{M}^4\times F$ 上で完全な標準模型のラグランジアンを生成する。
我々の枠組みはそのラグランジアンに現れる
\textit{結合定数の値}を生成する。

\subsection{導出されないもの}

式~\eqref{eq:alpha} の素数ギャップ $g(132)=5$ は
オイラー-マクローリン展開からは出ない。
これは追加の原理——物理的結合定数が $\Spec(\mathbb{Z})$ の
閉点（すなわち素数）に「量子化」されること——を必要とする。
構造的には自然だが（閉点はスキームの原子である）、予想の段階にとどまる。

因子 $4=d$（時空次元）も入力であり、導出されない。

\subsection{アインシュタイン (1905) との比較}

アインシュタインの論文は二つの公理（相対性原理 + 光速度不変）から
ローレンツ変換を導き、$E=mc^2$ を抽出した。
我々の論文は二つの公理（算術的時空 + スペクトル作用）から
ゼータ特殊値構造を導き、$\alpha^{-1}=137.036$ を抽出する。

特殊相対論と同様、我々の枠組みは\textbf{運動学的}である：
$\Spec(\mathbb{Z})$ 上の任意の力学理論が満たすべき拘束条件を決定するが、
完全な力学を指定しない。
このプログラムの「一般相対論」——アラケロフ幾何学による
$\Spec(\mathbb{Z})$ の曲率の決定——は今後の課題である。

%=====================================================================
\section{結論}
%=====================================================================

二つの公理から出発し、4つの厳密なステップで以下を導出した：

\begin{itemize}
\item $\Spec(\mathbb{Z})$ 上のスペクトル作用は
リーマンゼータ値の和である（ステップ1、厳密）。
\item その漸近展開はベルヌーイ数を必然的に含む
（ステップ2、標準数学）。
\item $\zeta(s)$ の $s=1$ での極が各項を
$\zeta(1{-}2j)$ に指数的に小さな補正で一致させる
（ステップ3、初等的複素解析）。
\item 結合定数 $12, 120, 252, 240, 132$ が
$2j/|B_{2j}|=1/|\zeta(1{-}2j)|$ として出現する
（ステップ4、ステップ1--3の合成）。
\end{itemize}

指数的に小さくない唯一の補正は $j=1$ の
$h'(0)=1-\frac{1}{2}\ln 2\pi$ であり、
これを $\Spec(\mathbb{Z})$ のアルキメデス的素点と同定する。

この枠組みが正しければ、微細構造定数は
自然の自由パラメータではなく、\textbf{整数の計算可能な不変量}である。

\bigskip

\begin{thebibliography}{99}
\bibitem{Einstein1905} A.~Einstein, Ann.\ Phys.\ \textbf{17}, 891 (1905).
\bibitem{BostConnes1995} J.-B.~Bost and A.~Connes, Selecta Math.\ \textbf{1}, 411 (1995).
\bibitem{Connes1996} A.~Connes, Commun.\ Math.\ Phys.\ \textbf{182}, 155 (1996).
\bibitem{ChamseddineConnes1997} A.~H.~Chamseddine and A.~Connes, Commun.\ Math.\ Phys.\ \textbf{186}, 731 (1997).
\bibitem{ConnesMarcolli2008} A.~Connes and M.~Marcolli, \textit{Noncommutative Geometry, Quantum Fields and Motives} (AMS, 2008).
\bibitem{Arakelov1974} S.~Arakelov, Math.\ USSR Izv.\ \textbf{8}, 1167 (1974).
\bibitem{Soule1992} C.~Soul\'e \textit{et al.}, \textit{Lectures on Arakelov Geometry} (Cambridge, 1992).
\bibitem{SSH1979} W.~P.~Su, J.~R.~Schrieffer, and A.~J.~Heeger, Phys.\ Rev.\ Lett.\ \textbf{42}, 1698 (1979).
\bibitem{Muong2_2021} B.~Abi \textit{et al.}\ (Muon $g{-}2$), Phys.\ Rev.\ Lett.\ \textbf{126}, 141801 (2021).
\end{thebibliography}

\end{document}
